\documentclass[12pt, a4paper]{report}

% =========================================================================
% PACKAGES
% =========================================================================
\usepackage{graphicx}
\usepackage{amsmath, amssymb}
\usepackage{setspace}
\usepackage{geometry}
\usepackage{fancyhdr}
\usepackage{titlesec}
\usepackage{tocloft}
\usepackage{longtable}
\usepackage{booktabs}
\usepackage{array}
\usepackage{float}
\usepackage{caption}
\usepackage{xcolor}
\usepackage{colortbl}
\usepackage{multirow}
\usepackage{enumitem}
\usepackage{listings}
\usepackage[most]{tcolorbox}
\usepackage{hyperref}

% =========================================================================
% PAGE SETUP & GEOMETRY
% =========================================================================
\geometry{margin=1in}
\onehalfspacing

% Colors matching AgenticPro & Academic Palette
\definecolor{primary}{HTML}{0F172A}      % Slate 900
\definecolor{brandteal}{HTML}{0F766E}    % Teal 700
\definecolor{teallight}{HTML}{CCFBF1}    % Teal 100
\definecolor{subtle}{HTML}{475569}       % Slate 600
\definecolor{bordercol}{HTML}{CBD5E1}    % Slate 300
\definecolor{bgalt}{HTML}{F8FAFC}        % Slate 50
\definecolor{successgreen}{HTML}{16A34A} % Green 600
\definecolor{skyaccent}{HTML}{0284C7}    % Sky 600

% Hyperref Setup
\hypersetup{
    colorlinks=true,
    linkcolor=brandteal,
    citecolor=brandteal,
    urlcolor=brandteal,
    pdftitle={Programming Lab Examination and Auto-Evaluation Agent},
    pdfauthor={Bhargavi K, Vaishnavi M, H B Samruddi, Adarsh K}
}

% Header & Footer Configuration
\pagestyle{fancy}
\fancyhf{}
\fancyhead[L]{\small\color{subtle}\textit{Programming Lab Examination \& Auto-Evaluation Agent}}
\fancyhead[R]{\small\color{subtle}\textit{Agentic AI}}
\fancyfoot[C]{\thepage}
\renewcommand{\headrulewidth}{0.4pt}
\renewcommand{\footrulewidth}{0.4pt}

% Chapter & Section Heading Styling
\titleformat{\chapter}[display]
  {\centering\normalfont\LARGE\bfseries\color{primary}}
  {\Large\bfseries\color{brandteal}Chapter \thechapter}
  {0.5em}
  {\Huge}
\titlespacing*{\chapter}{0pt}{-20pt}{25pt}

\titleformat{\section}
  {\Large\bfseries\color{primary}}
  {\thesection}{0.6em}{}
\titlespacing*{\section}{0pt}{14pt}{6pt}

\titleformat{\subsection}
  {\large\bfseries\color{brandteal}}
  {\thesubsection}{0.5em}{}
\titlespacing*{\subsection}{0pt}{10pt}{4pt}

% Table Column Helpers
\newcolumntype{L}[1]{>{\raggedright\arraybackslash}p{#1}}
\newcolumntype{C}[1]{>{\centering\arraybackslash}p{#1}}

% Code Listing Setup
\lstset{
    basicstyle=\ttfamily\footnotesize\color{primary},
    backgroundcolor=\color{bgalt},
    frame=single,
    rulecolor=\color{bordercol},
    breaklines=true,
    showstringspaces=false,
    tabsize=2
}

\begin{document}

% =========================================================================
% PAGE 1: TITLE & COVER PAGE
% =========================================================================
\begin{titlepage}
\thispagestyle{empty}
\centering

% University Logo / Header
\IfFileExists{kle logo.jpeg}{%
    \includegraphics[width=0.65\textwidth]{kle logo.jpeg}\\[0.2cm]
}{%
    \IfFileExists{kle_logo.jpeg}{%
        \includegraphics[width=0.65\textwidth]{kle_logo.jpeg}\\[0.2cm]
    }{%
        {\LARGE\bfseries\color{primary} KLE Technological University}\\[0.1cm]
        {\small\bfseries Dr. M. S. Sheshgiri Campus, Belagavi}\\[0.3cm]
    }%
}

\vspace{0.8cm}

{\huge\bfseries\color{primary} Programming Lab Examination and Auto-Evaluation Agent\par}
\vspace{0.4cm}
{\large\color{subtle}\textit{An Autonomous Multi-Agent Programming Assessment \& Evaluation Framework}\par}

\vspace{0.8cm}
{\normalsize\textbf{A Project Report}\par}
\vspace{0.2cm}
{\small Submitted in partial fulfillment of the requirements for the course\par}
\vspace{0.2cm}
{\large\bfseries\color{brandteal} Agentic AI\par}
\vspace{0.2cm}
{\small\textbf{Department of Computer Science and Engineering}\par}

\vspace{0.9cm}
{\small\textit{Submitted by:}\par}
\vspace{0.3cm}
\begin{tabular}{ll}
\textbf{Bhargavi K} & \textbf{02FE23BCI015} \\
\textbf{Vaishnavi M} & \textbf{02FE23BCI021} \\
\textbf{H B Samruddi} & \textbf{02FE23BCI052} \\
\textbf{Adarsh K} & \textbf{02FE23BCI014} \\
\end{tabular}

\vspace{0.9cm}
{\small\textit{Under the guidance of:}\par}
\vspace{0.2cm}
{\large\bfseries Prof. Priyanka Patil\par}
{\small Department of Computer Science and Engineering\par}
{\small KLE Technological University's Dr. M. S. Sheshgiri Campus, Belagavi\par}

\vfill
{\textbf{\normalsize KLE Technological University}\par}
{\small Dr. M. S. Sheshgiri Campus, Belagavi -- 590 008, Karnataka, India\par}
{\small\textbf{Academic Year 2026--2027}\par}

\end{titlepage}

% =========================================================================
% PAGE 2: DECLARATION
% =========================================================================
\newpage
\thispagestyle{plain}
\begin{center}
\IfFileExists{kle logo.jpeg}{%
    \includegraphics[width=0.45\textwidth]{kle logo.jpeg}\\[0.2cm]
}{%
    \IfFileExists{kle_logo.jpeg}{%
        \includegraphics[width=0.45\textwidth]{kle_logo.jpeg}\\[0.2cm]
    }{%
        {\large\bfseries KLE Technological University}\\[0.1cm]
    }%
}
\vspace{0.3cm}
{\LARGE\bfseries DECLARATION\par}
\end{center}

\vspace{0.5cm}
\noindent We hereby declare that the matter embodied in this report entitled \textbf{``Programming Lab Examination and Auto-Evaluation Agent''} submitted to \textbf{KLE Technological University} for the course completion of \textbf{Agentic AI} is the result of bona fide work carried out by us in the Department of Computer Science and Engineering, KLE Technological University's Dr. M. S. Sheshgiri Campus, Belagavi under the guidance of \textbf{Prof. Priyanka Patil}, Department of Computer Science and Engineering.

\vspace{0.3cm}
\noindent We further declare that to the best of our knowledge and belief, the work reported herein does not form part of any other project on the basis of which a course or award was conferred on an earlier occasion by any other student(s), and the results of the work have not been submitted for the award of any course, degree, or diploma within this or any other University or Institute. We hereby also confirm that all the software implementation, agent pipeline design, algorithmic problem curation, and automated verification tests presented in this report have been completed by us.

\vspace{1.2cm}
\noindent
\begin{tabular}{@{}l}
\textbf{Place:} Belagavi -- 590 008 \\
\textbf{Date:} 17 November 2026 \\
\end{tabular}

\vspace{1.5cm}
\noindent
\begin{tabular}{p{7.5cm}p{7.5cm}}
\textbf{Bhargavi K} & \textbf{H B Samruddi} \\
(02FE23BCI015) & (02FE23BCI052) \\[2.0cm]
\textbf{Vaishnavi M} & \textbf{Adarsh K} \\
(02FE23BCI021) & (02FE23BCI014) \\
\end{tabular}

\vfill
\noindent\rule{\textwidth}{0.4pt}
\centering{\footnotesize College of Engineering \& Technology, Udyambag, Belagavi -- 590 008, Karnataka, India.\\
Web: \url{www.klescet.ac.in}}

% =========================================================================
% PAGE 3: CERTIFICATE
% =========================================================================
\newpage
\thispagestyle{plain}
\begin{center}
\IfFileExists{kle logo.jpeg}{%
    \includegraphics[width=0.45\textwidth]{kle logo.jpeg}\\[0.2cm]
}{%
    \IfFileExists{kle_logo.jpeg}{%
        \includegraphics[width=0.45\textwidth]{kle_logo.jpeg}\\[0.2cm]
    }{%
        {\large\bfseries KLE Technological University}\\[0.1cm]
    }%
}
\vspace{0.3cm}
{\LARGE\bfseries CERTIFICATE\par}
\end{center}

\vspace{0.6cm}
\noindent This is to certify that the project entitled \textbf{``Programming Lab Examination and Auto-Evaluation Agent''} submitted to \textbf{KLE Technological University's Dr. M. S. Sheshgiri Campus, Belagavi} for the partial fulfillment of the requirements for the course \textbf{Agentic AI} is a bona fide record of work carried out by:
\begin{center}
\begin{tabular}{ll}
\textbf{Bhargavi K} & (02FE23BCI015) \\
\textbf{Vaishnavi M} & (02FE23BCI021) \\
\textbf{H B Samruddi} & (02FE23BCI052) \\
\textbf{Adarsh K} & (02FE23BCI014) \\
\end{tabular}
\end{center}
\noindent students in the Department of Computer Science and Engineering, KLE Technological University's Dr. M. S. Sheshgiri Campus, Belagavi, under my supervision and guidance. The contents of this report, in full or in parts, have not been submitted to any other Institute or University for the award of any degree, diploma, or course credit.

\vspace{2.0cm}
\noindent
\begin{tabular}{@{}l}
\textbf{Place:} Belagavi -- 590 008 \\
\textbf{Date:} 17 November 2026 \\
\end{tabular}

\vspace{2.2cm}
\noindent
\begin{tabular}{p{7.5cm}p{7.5cm}}
\textbf{Prof. Priyanka Patil} & \textbf{Dr. V. S. Rajpurohit} \\
Course Guide \& Coordinator & Professor \& Head of the Department \\
Dept. of Computer Science \& Engineering & Dept. of Computer Science \& Engineering \\
KLE Tech University, Dr. MSSCET & KLE Tech University, Dr. MSSCET \\
\end{tabular}

\vfill
\noindent\rule{\textwidth}{0.4pt}
\centering{\footnotesize College of Engineering \& Technology, Udyambag, Belagavi -- 590 008, Karnataka, India.\\
Web: \url{www.klescet.ac.in}}

% =========================================================================
% PAGE 4: ABSTRACT & SYSTEM OVERVIEW
% =========================================================================
\newpage
\pagenumbering{roman}
\setcounter{page}{4}

\begin{center}
{\LARGE\bfseries Abstract\par}
\end{center}
\vspace{0.5cm}

\noindent Traditional automated grading platforms (such as LeetCode, HackerRank, and DOMjudge) rely almost exclusively on binary input/output assertion tests executed within sandboxed processes. While adequate for verifying input-output correctness, this approach fails to assess software engineering fundamentals: code quality, cyclomatic complexity, algorithmic asymptotic efficiency, runtime failure root causes, security blast radius, AST semantic grounding, and domain constraints.

\vspace{0.3cm}
\noindent To solve these challenges, this project presents the \textbf{Programming Lab Examination and Auto-Evaluation Agent} (built on the \textbf{AgenticPro} platform), an autonomous, multi-agent evaluation suite specifically tailored for university-level computer science laboratories. The system incorporates a multi-agent swarm across 14 laboratory disciplines:
\begin{itemize}[itemsep=2pt,leftmargin=1.5em]
    \item \textbf{Labs 1 to 8 (Foundational Evaluator):} Secure process sandboxing, multi-case assertion testing, AST syntax inspection, Big-$\mathcal{O}$ empirical curve fitting, security taint detection, pedagogical hint synthesis, and weighted rubric grading.
    \item \textbf{Labs 9 to 14 (Systemic Engineering Capabilities):} Deep literature research, cloud incident root-cause triage (OOMKilled, Deadlocks), autonomous SDLC code generation, pre-execution safety gating, engineering physics validation, and team capstone collaboration.
\end{itemize}

\vspace{0.3cm}
\noindent A core design principle is the \textbf{clean separation of concerns}: Labs 9--14 run behind the scenes to enrich the student evaluation panel rather than cluttering the UI with disconnected playground tabs. When a student submits a program, the system returns test case outputs alongside a Lab 12 Safety Risk Badge, Lab 10 Runtime Incident Diagnostics, Lab 11 SDLC Quality Grade, Lab 9 Semantic Grounding Audit, and Lab 14 Multi-Agent Consensus Scoring. The platform includes a curated \textbf{27-problem pure Data Structures and Algorithms (DSA) Question Bank} spanning 8 domains (Arrays, Strings, Stacks, Queues, Linked Lists, Trees, Graphs, and Dynamic Programming) across three difficulty tiers (L1, L2, L3) with real-world scenarios.

\vspace{0.4cm}
% Summary Table
\begin{center}
\renewcommand{\arraystretch}{1.2}
\small
\begin{tabular}{|L{3.5cm}|L{11.5cm}|}
\hline
\rowcolor{brandteal}
\color{white}\textbf{Metric / Component} & \color{white}\textbf{Implementation Specification} \\
\hline
\textbf{Course \& University} & Agentic AI $\cdot$ KLE Technological University, Dr. MSSCET, Belagavi \\
\hline
\rowcolor{bgalt}
\textbf{Project Title} & Programming Lab Examination and Auto-Evaluation Agent \\
\hline
\textbf{Technology Stack} & FastAPI (Python 3.10), React 19, SQLite, Vite, Pydantic, ReportLab \\
\hline
\rowcolor{bgalt}
\textbf{Verification Suite} & \textbf{\color{successgreen}49 / 49 Unit \& Integration Tests Passed (100\% Pass Rate, 3.85s)} \\
\hline
\textbf{Curriculum Scope} & 14 Laboratories $\cdot$ 27 Pure-DSA Scenario Problems $\cdot$ 8 Domains $\cdot$ L1--L3 Tiers \\
\hline
\rowcolor{bgalt}
\textbf{Resilience Model} & Dual-Path Fallback (0.5s pre-flight probe with zero-downtime offline synthesis) \\
\hline
\end{tabular}
\end{center}

% =========================================================================
% PAGE 5: TABLE OF CONTENTS & LIST OF TABLES
% =========================================================================
\newpage
\tableofcontents

\vspace{1.0cm}
\listoftables

\newpage
\pagenumbering{arabic}
\setcounter{page}{1}

% =========================================================================
% CHAPTER 1: INTRODUCTION & EXECUTIVE OVERVIEW
% =========================================================================
\chapter{Introduction \& Executive Overview}

\section{Project Motivation}
Programming laboratory examinations in undergraduate computer science curricula present substantial pedagogical and logistical challenges. Faculty members must continuously curate original questions, author edge-case test suites, prevent plagiarism, and provide constructive feedback across hundreds of student submissions under tight deadlines. 

Existing automated judges provide instantaneous execution results but lack the cognitive depth of human evaluators:
\begin{enumerate}[itemsep=2pt,leftmargin=1.5em]
    \item \textbf{Binary Pass/Fail Limitation:} A single off-by-one error or whitespace discrepancy results in zero points, failing to assess algorithmic structure or partial correctness.
    \item \textbf{Absence of Security \& Safety Gating:} Untrusted student code is frequently executed without rigorous pre-execution AST taint analysis, risking host container compromise or Denial of Service (DoS).
    \item \textbf{Zero Root-Cause Triage:} When code encounters a runtime exception (e.g., Memory Limit Exceeded or Infinite Recursion), standard judges provide opaque error strings rather than actionable diagnostic remediation.
    \item \textbf{Superficial Metric Analysis:} Code quality, maintainability, cyclomatic complexity, and documentation compliance are ignored.
\end{enumerate}

\section{The AgenticPro Solution}
The \textbf{Programming Lab Examination and Auto-Evaluation Agent} bridges this gap by introducing a coordinated multi-agent artificial intelligence swarm that evaluates code as an integrated software engineering artifact. By combining deterministic execution sandboxes with specialized LLM agent personas (Lead Architect, QA Lead, CloudOps Engineer, Governance Officer), the system provides instant, comprehensive, and pedagogically rich assessments.

\begin{center}
\renewcommand{\arraystretch}{1.3}
\begin{tabular}{|C{3.0cm}|C{3.0cm}|C{3.0cm}|C{3.0cm}|C{3.0cm}|}
\hline
\rowcolor{bgalt}
\textbf{14 Labs} & \textbf{49 Tests} & \textbf{27 DSA Problems} & \textbf{14 Tables} & \textbf{Zero-Downtime} \\
\rowcolor{bgalt}
{\footnotesize\color{subtle}Fully Implemented} & 
{\footnotesize\color{successgreen}100\% Passed (3.85s)} & 
{\footnotesize\color{skyaccent}8 Topics (L1--L3)} & 
{\footnotesize\color{subtle}Relational Schema} & 
{\footnotesize\color{brandteal}Offline Synthesizer} \\
\hline
\end{tabular}
\end{center}

% =========================================================================
% CHAPTER 2: SYSTEM ARCHITECTURE & TOPOLOGY
% =========================================================================
\chapter{System Architecture \& Multi-Tier Topology}

\section{Architectural Layers}
The platform is engineered using a decoupled, asynchronous multi-tier architecture to ensure high throughput, deterministic grading, and resilient fault isolation:

\begin{table}[H]
\centering
\caption{System Architecture Tier Breakdown}
\renewcommand{\arraystretch}{1.2}
\small
\begin{tabular}{|L{2.5cm}|L{3.2cm}|L{3.0cm}|L{6.8cm}|}
\hline
\rowcolor{brandteal}
\color{white}\textbf{Tier} & \color{white}\textbf{Component Name} & \color{white}\textbf{Technologies} & \color{white}\textbf{Architectural Responsibility} \\
\hline
\textbf{Presentation} & Role-Guarded Web UI & React 19, Vite, CSS Tokens & Student code arena, real-time multi-agent results panel, faculty exam creator, and analytics gradebook. \\
\hline
\rowcolor{bgalt}
\textbf{API Gateway} & FastAPI REST Router & FastAPI, Starlette, CORS & 12 router modules exposing endpoints for exams, submissions, evaluations, and individual lab agents. \\
\hline
\textbf{Agent Swarm} & Evaluation Swarm Engine & Pydantic, Ollama Adapter, Fallback Synthesizer & Orchestrates Lead Architect, QA Lead, CloudOps Triage, and Governance Officer in a multi-agent consensus pipeline. \\
\hline
\rowcolor{bgalt}
\textbf{Execution Sandbox} & Isolated Sandbox Worker & Subprocess, OS Jail, AST & Executes untrusted student code within CPU/memory quotas (5.0s timeout, restricted syscalls). \\
\hline
\textbf{Persistence} & SQLite ACID Store & SQLite3, WAL Mode & Stores users, exams, questions, test cases, submissions, evaluations, and domain-specific audit records. \\
\hline
\end{tabular}
\end{table}

\section{Dual-Path Heuristic Fallback Principle}
In educational deployments, platform availability is paramount. Network instability or LLM server congestion must never impede students from submitting exam solutions. 

To ensure zero downtime, all agents implement a \textbf{Dual-Path Resilient Pipeline}:
\begin{enumerate}[itemsep=2pt,leftmargin=1.5em]
    \item \textbf{0.5-Second Pre-Flight Probe:} Prior to dispatching an evaluation payload, the backend issues an asynchronous health probe to the active LLM provider (Ollama / Gemini endpoint).
    \item \textbf{Dynamic Failover:} If the response time exceeds 500ms or returns an HTTP error, the system dynamically routes the request to high-fidelity, deterministic domain synthesizers.
    \item \textbf{Guaranteed SLA:} Student evaluations, test case executions, and exam creation operations complete with 100\% reliability and zero hung connections.
\end{enumerate}

% =========================================================================
% CHAPTER 3: CURRICULUM BREAKDOWN: LABS 1 TO 8
% =========================================================================
\chapter{Foundational Curriculum: Labs 1 to 8 (Core Evaluator)}

\section{Overview}
Labs 1 through 8 form the foundational execution and evaluation core of the platform. These units establish deterministic code execution, static syntax validation, algorithmic complexity estimation, security sandboxing, and multi-criteria scoring.

\begin{table}[H]
\centering
\caption{Labs 1 to 8: Foundational Evaluation Engine}
\renewcommand{\arraystretch}{1.2}
\small
\begin{tabular}{|L{1.5cm}|L{3.8cm}|L{5.0cm}|L{5.2cm}|}
\hline
\rowcolor{brandteal}
\color{white}\textbf{Lab} & \color{white}\textbf{Title \& Scope} & \color{white}\textbf{Core Capabilities} & \color{white}\textbf{Verification Oracle / Criteria} \\
\hline
\textbf{Lab 1} & Automated Unit Testing \& IO & Standard I/O stream piping, timeout handling, exit code inspection. & Exact output matching, whitespace normalization, returncode == 0. \\
\hline
\rowcolor{bgalt}
\textbf{Lab 2} & AST Code Quality \& Linting & Static Python AST inspection, PEP8 style conformance, naming conventions. & Flake8 violation count, cyclomatic complexity index. \\
\hline
\textbf{Lab 3} & Algorithmic Complexity \& Big-$\mathcal{O}$ & Dynamic profiling across multiple input sizes ($N=10$ to $N=100{,}000$). & Empirical asymptotic curve fitting ($\mathcal{O}(1), \mathcal{O}(N), \mathcal{O}(N \log N), \mathcal{O}(N^2)$). \\
\hline
\rowcolor{bgalt}
\textbf{Lab 4} & Security \& Taint Analysis & Screening for dangerous calls (\texttt{os.system}, \texttt{eval}, \texttt{open}, \texttt{subprocess}). & AST node security visitor, malicious process containment. \\
\hline
\textbf{Lab 5} & Semantic AI Feedback & Pedagogical hint generation without revealing outright solutions. & Socratic guidance rubric, tone compliance, learning hint clarity. \\
\hline
\rowcolor{bgalt}
\textbf{Lab 6} & Dynamic Process Sandboxing & CPU time throttling (5s max), memory limits, file descriptor restrictions. & SIGKILL enforcement on runaway loops, isolation boundary checks. \\
\hline
\textbf{Lab 7} & Rubric-Based Grading & Multi-criteria scoring: Correctness (50\%), Quality (20\%), Efficiency (30\%). & Mathematical weighted formula, point scaling to question max. \\
\hline
\rowcolor{bgalt}
\textbf{Lab 8} & Swarm Orchestration Core & Multi-agent message passing between Evaluator, Quality, and Security agents. & Synthesized consensus score, conflict resolution protocol. \\
\hline
\end{tabular}
\end{table}

% =========================================================================
% CHAPTER 4: ADVANCED CURRICULUM: LABS 9 TO 14
% =========================================================================
\chapter{Advanced Systemic Curriculum: Labs 9 to 14}

\section{Overview \& Systemic Expansion}
In modern software engineering, evaluating code extends far beyond isolated functions. Labs 9 through 14 elevate the student experience to evaluate the \textbf{entire autonomous software engineering lifecycle}: verifying external API documentation without hallucinations, diagnosing runtime cloud crashes, synthesizing full codebases from specifications, auditing security blast radius, verifying physical engineering laws, and running full multi-agent capstone studios.

\begin{table}[H]
\centering
\caption{Labs 9 to 14: Systemic Engineering Capabilities}
\renewcommand{\arraystretch}{1.2}
\small
\begin{tabular}{|L{1.5cm}|L{3.5cm}|L{5.2cm}|L{5.3cm}|}
\hline
\rowcolor{brandteal}
\color{white}\textbf{Lab} & \color{white}\textbf{Specialized Agent} & \color{white}\textbf{Domain \& Problem Solved} & \color{white}\textbf{Evaluation Output \& Metrics} \\
\hline
\textbf{Lab 9} & Deep Research Agent & Pre-coding literature research, API documentation verification, and anti-hallucination grounding. & Source credibility scores, citation verification matrix, and grounding verdict. \\
\hline
\rowcolor{bgalt}
\textbf{Lab 10} & Cloud Operations Agent & Automated post-mortem diagnosis for runtime crashes (OOMKilled, CrashLoopBackOff, SSL expiry, DB pools). & Root-cause categorization, Mean Time to Diagnosis (MTTD), and automated remediation script. \\
\hline
\textbf{Lab 11} & Agentic SDLC Factory & Autonomous specification-to-code synthesis across 3 tiers: Architecture $\rightarrow$ Source $\rightarrow$ Pytest. & Modularity score, cyclomatic complexity index, test harness pass rate (100\%). \\
\hline
\rowcolor{bgalt}
\textbf{Lab 12} & Safety \& Governance Agent & Pre-execution safety gating, risk tiering (LOW, MED, HIGH, CRITICAL), and blast radius analysis. & Blast radius score (0--100), dual-key signoff gate, and automated rollback plan. \\
\hline
\textbf{Lab 13} & Discipline-Specific Agents & Engineering physics verification across Mechanical (ME), Civil (CE), and Electrical (EE) standards. & Safety factors ($\ge 1.5$), physical conservation invariants, and code compliance verdicts. \\
\hline
\rowcolor{bgalt}
\textbf{Lab 14} & Capstone Project Studio & Full team-owned agentic system: 5-member swarm (Architect, Researcher, Engineer, QA, Governance) in 5 stages. & Deliverable explorer, in-memory test harness, zip packaging, and markdown audit dossier. \\
\hline
\end{tabular}
\end{table}

\section{Integration into the Live Program Evaluation UI}
A crucial architectural decision in this project was \textbf{not to clutter the user interface with isolated, theoretical playground pages for Labs 9--14}. Instead, the capabilities developed in Labs 9--14 run directly on the backend and power the \textbf{Student Evaluation Result Panel}:

\begin{itemize}[itemsep=4pt,leftmargin=1.5em]
    \item \textbf{Lab 12 Pre-Execution Safety Gate:} Every student submission is passed through an AST safety visitor before execution. If prohibited system calls (e.g., \texttt{import os; os.remove}) are detected, execution is halted immediately and a \textbf{CRITICAL Risk} badge with a blast radius breakdown is displayed. Safe code receives a \textbf{LOW Risk (Approved)} badge.
    \item \textbf{Lab 10 Runtime Crash Diagnostics:} If a student's code fails test cases due to an infinite loop (Timeout), high memory consumption, or recursion depth errors, the CloudOps triage agent automatically diagnoses the root cause and provides a concrete code fix in the results panel.
    \item \textbf{Lab 11 SDLC Quality \& Complexity Grade:} The system parses student code into an Abstract Syntax Tree to compute cyclomatic complexity, function nesting depth, and adherence to PEP8, producing an objective \textbf{Clean Code Grade (A/B/C)}.
    \item \textbf{Lab 9 Semantic Grounding Check:} The submission is verified to ensure all imported packages are part of approved standard libraries, preventing hallucinated or unsupported dependencies.
    \item \textbf{Lab 14 Swarm Committee Consensus:} Evaluation scores from the Lead Architect, QA Lead, and Governance Officer are aggregated into a weighted consensus score with multi-perspective qualitative feedback.
\end{itemize}

% =========================================================================
% CHAPTER 5: PURE-DSA QUESTION BANK & DIFFICULTY HIERARCHY
% =========================================================================
\chapter{Pure-DSA Question Bank \& Difficulty Hierarchy}

\section{Curriculum Design Principles}
To ensure academic rigor and domain focus, the Question Bank was completely restructured around core \textbf{Data Structures and Algorithms (DSA)}. Every question is framed within a realistic engineering scenario rather than an abstract mathematical riddle, requiring students to apply data structure concepts to solve practical systems problems.

Questions are categorized into three standardized difficulty tiers:
\begin{table}[H]
\centering
\caption{Difficulty Tier Definitions}
\renewcommand{\arraystretch}{1.2}
\small
\begin{tabular}{|L{3.0cm}|L{4.2cm}|L{8.3cm}|}
\hline
\rowcolor{brandteal}
\color{white}\textbf{Tier} & \color{white}\textbf{Target Skill Level} & \color{white}\textbf{Algorithmic \& Data Structure Scope} \\
\hline
\textbf{L1} (Foundational) & Core Syntax \& Linear Scans & Basic indexing, single loops, condition checks, array scans, and direct string/queue manipulations. \\
\hline
\rowcolor{bgalt}
\textbf{L2} (Intermediate) & Multi-Step Data Structures & Monotonic stacks, circular queues, two-pointer windowing, tree BFS/DFS traversals, and Floyd cycle detection. \\
\hline
\textbf{L3} (Advanced) & Complex Algorithms \& Optimization & Dynamic programming, Dijkstra shortest paths, lowest common ancestors, monotonic deques, and edit distances. \\
\hline
\end{tabular}
\end{table}

\section{Complete Scenario-Based DSA Question Catalog}
The curriculum comprises \textbf{27 questions across 8 core DSA domains}, fully seeded in the relational database with 54 automated input/output test cases:

\begin{center}
\renewcommand{\arraystretch}{1.15}
\small
\begin{longtable}{|L{2.2cm}|L{2.0cm}|C{0.9cm}|L{3.8cm}|C{0.9cm}|L{6.2cm}|}
\caption{Complete 27-Question Pure-DSA Scenario Catalog} \label{tab:dsa_catalog} \\
\hline
\rowcolor{brandteal}
\color{white}\textbf{ID} & \color{white}\textbf{Domain} & \color{white}\textbf{Tier} & \color{white}\textbf{Scenario Title} & \color{white}\textbf{Pts} & \color{white}\textbf{Real-World Scenario \& DSA Concept} \\
\hline
\endfirsthead

\hline
\rowcolor{brandteal}
\color{white}\textbf{ID} & \color{white}\textbf{Domain} & \color{white}\textbf{Tier} & \color{white}\textbf{Scenario Title} & \color{white}\textbf{Pts} & \color{white}\textbf{Real-World Scenario \& DSA Concept} \\
\hline
\endhead

\hline
\multicolumn{6}{|r|}{\footnotesize\color{subtle}Continued on next page\dots} \\
\hline
\endfoot

\hline
\endlastfoot

% Arrays
\texttt{q-test-001} & Arrays & \textbf{L1} & Cashier Checkout Subtotal & 10 & Barcode dual-price parsing \& summation \\
\hline
\rowcolor{bgalt}
\texttt{q-test-002} & Arrays & \textbf{L1} & Server Peak CPU Spike Detector & 15 & Cloud telemetry array maximum detection \\
\hline
\texttt{q-dsa-arr-l1} & Arrays & \textbf{L1} & E-Commerce Voucher Pairing & 5 & Cart item complement hash index lookup \\
\hline
\rowcolor{bgalt}
\texttt{q-dsa-arr-l2} & Arrays & \textbf{L2} & Smart Meter Power Window & 10 & Grid telemetry prefix sum target window \\
\hline
\texttt{q-dsa-arr-l3} & Arrays & \textbf{L3} & Skyline Stormwater Retention & 15 & Skyscraper rainwater trapping $\mathcal{O}(N)$ \\
\hline

% Strings
\rowcolor{bgalt}
\texttt{q-test-003} & Strings & \textbf{L1} & DNA Palindromic Symmetry & 10 & Genomic nucleotide string symmetry validator \\
\hline
\texttt{q-dsa-str-l1} & Strings & \textbf{L1} & Cryptographic Anagram Auth & 5 & Access credential frequency matching \\
\hline
\rowcolor{bgalt}
\texttt{q-dsa-str-l2} & Strings & \textbf{L2} & Signal Unique Packet Window & 10 & 5G radio non-repeating signal window \\
\hline
\texttt{q-dsa-str-l3} & Strings & \textbf{L3} & Genomic Mutation Distance & 15 & RNA mutation Levenshtein edit distance \\
\hline

% Stacks
\rowcolor{bgalt}
\texttt{q-dsa-stk-l1} & Stacks & \textbf{L1} & Compiler Delimiter Validator & 5 & Code IDE syntax parser scope balancing \\
\hline
\texttt{q-dsa-stk-l2} & Stacks & \textbf{L2} & Stock Next Breakout Day & 10 & Trading desk monotonic price alert stack \\
\hline
\rowcolor{bgalt}
\texttt{q-dsa-stk-l3} & Stacks & \textbf{L3} & Skyline Billboard Optimization & 15 & Urban LED billboard histogram max area \\
\hline

% Queues
\texttt{q-dsa-que-l1} & Queues & \textbf{L1} & Cloud Print Spooler Dispatch & 5 & Enterprise FIFO print task spooling \\
\hline
\rowcolor{bgalt}
\texttt{q-dsa-que-l2} & Queues & \textbf{L2} & Chat Unique Character Dispatch & 10 & Live support stream first non-repeat \\
\hline
\texttt{q-dsa-que-l3} & Queues & \textbf{L3} & Router Peak Traffic Sliding Window & 15 & Network packet burst monotonic deque \\
\hline

% Linked Lists
\rowcolor{bgalt}
\texttt{q-dsa-ll-l1} & Linked Lists & \textbf{L1} & Browser History Reversal & 5 & Navigation trail pointer in-place reversal \\
\hline
\texttt{q-dsa-ll-l2} & Linked Lists & \textbf{L2} & Supply Chain Inventory Merge & 10 & Automated logistics sorted feed splice \\
\hline
\rowcolor{bgalt}
\texttt{q-dsa-ll-l3} & Linked Lists & \textbf{L3} & RPC Call Chain Loop Detection & 15 & Microservice deadlock cycle detection \\
\hline

% Trees
\texttt{q-dsa-tree-l1} & Trees & \textbf{L1} & Org Chart Management Depth & 5 & Corporate managerial tree hierarchy depth \\
\hline
\rowcolor{bgalt}
\texttt{q-dsa-tree-l2} & Trees & \textbf{L2} & Warehouse Drone Shelf Tier Scan & 10 & Vertical rack BFS level order traversal \\
\hline
\texttt{q-dsa-tree-l3} & Trees & \textbf{L3} & DNS Lowest Common Authority & 15 & Distributed naming tree lowest ancestor \\
\hline

% Graphs
\rowcolor{bgalt}
\texttt{q-dsa-graph-l1} & Graphs & \textbf{L1} & Smart City Hub Direct Degree & 5 & Municipal optical grid junction degrees \\
\hline
\texttt{q-dsa-graph-l2} & Graphs & \textbf{L2} & Social Isolated Communities & 10 & Campus social network disjoint clusters \\
\hline
\rowcolor{bgalt}
\texttt{q-dsa-graph-l3} & Graphs & \textbf{L3} & Ambulance Emergency GPS Route & 15 & 911 road network Dijkstra shortest path \\
\hline

% Dynamic Programming
\texttt{q-dsa-dp-l1} & DP & \textbf{L1} & Mars Rover Motor Stride Modes & 5 & Robotic rover climbing stair recurrence \\
\hline
\rowcolor{bgalt}
\texttt{q-dsa-dp-l2} & DP & \textbf{L2} & ATM Cash Dispenser Optimization & 10 & Vault bill transfer coin change DP \\
\hline
\texttt{q-dsa-dp-l3} & DP & \textbf{L3} & Rocket Payload Value Maximization & 15 & Satellite mission 0/1 knapsack weight \\
\hline
\end{longtable}
\end{center}

% =========================================================================
% CHAPTER 6: RELATIONAL DATABASE SCHEMA & DATA MODELS
% =========================================================================
\chapter{Relational Database Schema \& Data Models}

The application utilizes an ACID-compliant relational SQLite store (\texttt{evaluations.db}) with Write-Ahead Logging (WAL) mode enabled for concurrent reads and writes:

\begin{table}[H]
\centering
\caption{Database Schema Table Directory}
\renewcommand{\arraystretch}{1.2}
\small
\begin{tabular}{|L{3.2cm}|L{2.5cm}|L{6.8cm}|L{3.2cm}|}
\hline
\rowcolor{brandteal}
\color{white}\textbf{Table Name} & \color{white}\textbf{Primary Key} & \color{white}\textbf{Key Fields \& Data Stored} & \color{white}\textbf{Associated Module} \\
\hline
\texttt{users} & \texttt{id} & username, password\_hash, role (student/faculty), full\_name, email & Auth System \\
\hline
\rowcolor{bgalt}
\texttt{exams} & \texttt{id} & title, starts\_at, duration\_minutes & Exam Management \\
\hline
\texttt{questions} & \texttt{id} & title, prompt, language, points, difficulty (L1--L3), domain & Question Bank \\
\hline
\rowcolor{bgalt}
\texttt{exam\_questions} & \texttt{(exam\_id, question\_id)} & Foreign keys linking exams to associated questions & Exam Assembler \\
\hline
\texttt{test\_cases} & \texttt{id} & question\_id, stdin, expected\_stdout, points, is\_hidden & Test Runner \\
\hline
\rowcolor{bgalt}
\texttt{submissions} & \texttt{id} & exam\_id, question\_id, student\_id, language, source\_code, submitted\_at & Student Arena \\
\hline
\texttt{evaluations} & \texttt{submission\_id} & passed, total, score, feedback, status, \textbf{agent\_insights} (JSON) & Multi-Agent Engine \\
\hline
\rowcolor{bgalt}
\texttt{results} & \texttt{(student\_id, exam\_id)} & score, max\_score, published (boolean) & Official Gradebook \\
\hline
\texttt{research\_sessions} & \texttt{id} & topic, synthesized\_report, sources\_json, created\_at & Lab 9 (Research) \\
\hline
\rowcolor{bgalt}
\texttt{cloud\_ops\_incidents} & \texttt{id} & service\_name, severity, root\_cause, report\_json, status & Lab 10 (CloudOps) \\
\hline
\texttt{sdlc\_projects} & \texttt{id} & project\_name, requirement\_text, artifact\_json, status & Lab 11 (SDLC) \\
\hline
\rowcolor{bgalt}
\texttt{governance\_audits} & \texttt{id} & action\_title, overall\_risk\_score, risk\_tier, status, audit\_json & Lab 12 (Governance) \\
\hline
\texttt{discipline\_worklets} & \texttt{id} & discipline (ME/CE/EE), compliance\_verdict, worklet\_json & Lab 13 (Discipline) \\
\hline
\rowcolor{bgalt}
\texttt{capstone\_projects} & \texttt{id} & team\_name, project\_title, active\_stage, deliverables\_json, test\_results\_json & Lab 14 (Capstone) \\
\hline
\end{tabular}
\end{table}

% =========================================================================
% CHAPTER 7: QUALITY ASSURANCE & TEST VERIFICATION REPORT
% =========================================================================
\chapter{Quality Assurance \& Test Verification Report}

Every endpoint, database operation, AST security check, and multi-agent synthesis pipeline is covered by automated unit and integration tests using \texttt{pytest}:

\begin{table}[H]
\centering
\caption{Automated Test Verification Matrix}
\renewcommand{\arraystretch}{1.2}
\small
\begin{tabular}{|L{4.5cm}|L{3.2cm}|L{6.0cm}|C{2.3cm}|}
\hline
\rowcolor{brandteal}
\color{white}\textbf{Test Module Path} & \color{white}\textbf{Target Component} & \color{white}\textbf{Test Scope \& Assertions} & \color{white}\textbf{Status} \\
\hline
\texttt{tests/test\_capstone\_studio.py} & Lab 14: Capstone Studio & Models, in-memory zip packaging, 5-stage pipeline, SQLite persistence, API routes & \textbf{PASSED} (6/6) \\
\hline
\rowcolor{bgalt}
\texttt{tests/test\_discipline\_agents.py} & Lab 13: Discipline Agents & ME thermal, CE beam shear, EE power flow, compliance verification & \textbf{PASSED} (6/6) \\
\hline
\texttt{tests/test\_safety\_governance.py} & Lab 12: Safety Governance & Risk score calculation, blast radius, dual-key authorization gate, rollback scripts & \textbf{PASSED} (4/4) \\
\hline
\rowcolor{bgalt}
\texttt{tests/test\_sdlc\_factory.py} & Lab 11: SDLC Factory & Requirement parsing, architecture synthesis, code generator, test harness runner & \textbf{PASSED} (5/5) \\
\hline
\texttt{tests/test\_cloud\_ops.py} & Lab 10: Cloud Operations & OOMKilled, CrashLoopBackOff, SSL expiry, DB connection pool remediation & \textbf{PASSED} (4/4) \\
\hline
\rowcolor{bgalt}
\texttt{tests/test\_deep\_research.py} & Lab 9: Deep Research & Search triangulation, citation verification, summary synthesis, fallback safety & \textbf{PASSED} (4/4) \\
\hline
\texttt{tests/api/*.py} & Core API Endpoints & Auth, exams, questions, submissions, test cases, and evaluation routes & \textbf{PASSED} (12/12) \\
\hline
\rowcolor{bgalt}
\texttt{tests/skills/*.py} & Core Evaluation Skills & Code execution sandbox, AST security parsing, rubric scoring, swarm adapter & \textbf{PASSED} (8/8) \\
\hline
\rowcolor{teallight}
\textbf{TOTAL SUITE RESULT} & \textbf{Full Assessment Suite} & \textbf{49 unit \& integration tests executed across all 14 laboratories} & \textbf{\color{successgreen}49 PASSED (3.85s)} \\
\hline
\end{tabular}
\end{table}

% =========================================================================
% CHAPTER 8: DEPLOYMENT & OPERATOR'S RUNBOOK
% =========================================================================
\chapter{Deployment \& Operator's Runbook}

\section{System Prerequisites}
\begin{itemize}[itemsep=2pt,leftmargin=1.5em]
    \item Python 3.10 or higher with FastAPI, Uvicorn, Pydantic, and SQLite3.
    \item Node.js 18+ and npm for the React 19 / Vite client.
\end{itemize}

\section{Backend Service Startup}
In Terminal 1:
\begin{lstlisting}[language=bash]
cd AgenticPro
python -m uvicorn backend.main:app --reload --port 8000
\end{lstlisting}
\begin{itemize}[itemsep=1pt,leftmargin=1.5em]
    \item REST API Base URL: \url{http://localhost:8000}
    \item Interactive Swagger OpenAPI Documentation: \url{http://localhost:8000/docs}
    \item Health Check Endpoint: \url{http://localhost:8000/api/stats}
\end{itemize}

\section{Frontend Client Startup}
In Terminal 2:
\begin{lstlisting}[language=bash]
cd AgenticPro/frontend
npm run dev
\end{lstlisting}
\begin{itemize}[itemsep=1pt,leftmargin=1.5em]
    \item Web UI URL: \url{http://localhost:5173}
    \item Production Build Command: \texttt{npm run build} (built in 3.00s with 0 errors)
\end{itemize}

\section{Default User Credentials}
\begin{table}[H]
\centering
\caption{Default User Roles and Authentication}
\renewcommand{\arraystretch}{1.2}
\small
\begin{tabular}{|L{3.5cm}|L{2.5cm}|L{2.5cm}|L{7.5cm}|}
\hline
\rowcolor{brandteal}
\color{white}\textbf{Portal Role} & \color{white}\textbf{Username} & \color{white}\textbf{Password} & \color{white}\textbf{Permitted Capabilities} \\
\hline
\textbf{Student Workspace} & \texttt{student1} & \texttt{student123} & View active exams, code in sandbox arena, submit solutions, inspect real-time multi-agent evaluations. \\
\hline
\rowcolor{bgalt}
\textbf{Faculty Workspace} & \texttt{faculty1} & \texttt{faculty123} & Curate question bank, associate questions by domain/level, create exams, view gradebook, publish results. \\
\hline
\end{tabular}
\end{table}

% =========================================================================
% CHAPTER 9: CONCLUSION & PROJECT SIGN-OFF
% =========================================================================
\chapter{Conclusion \& Project Sign-Off}

\section{Project Achievements}
The \textbf{Programming Lab Examination and Auto-Evaluation Agent} developed for the \textbf{Agentic AI} course successfully demonstrates how multi-agent LLM systems can transform programming education. Key achievements include:
\begin{enumerate}[itemsep=2pt,leftmargin=1.5em]
    \item \textbf{Comprehensive Evaluation:} Seamless multi-agent assessment covering unit tests, AST syntax, Big-$\mathcal{O}$ complexity, security taint analysis, and runtime diagnostics.
    \item \textbf{Production UI Architecture:} Clean separation of concerns with student and faculty workflows front-and-center, running systemic agents behind the scenes.
    \item \textbf{Curriculum Rigor:} 27 pure-DSA scenario-based questions across 8 domains and L1--L3 tiers.
    \item \textbf{Zero-Downtime Resilience:} Dual-path 0.5s pre-flight fallback preventing platform stalls.
    \item \textbf{100\% Verified Codebase:} 49/49 passing automated pytest test cases and clean React 19 production builds.
\end{enumerate}

\vspace{1.0cm}
\section{Team Sign-Off}
\begin{table}[H]
\centering
\caption{Project Architectural Sign-Off}
\renewcommand{\arraystretch}{1.5}
\small
\begin{tabular}{|L{7.5cm}|L{7.5cm}|}
\hline
\rowcolor{bgalt}
\textbf{Student Development Team:} & \textbf{Course Guide \& Approval:} \\
Bhargavi K (02FE23BCI015) & Prof. Priyanka Patil \\
Vaishnavi M (02FE23BCI021) & Course Guide \& Assistant Professor \\
H B Samruddi (02FE23BCI052) & Department of Computer Science \& Engineering \\
Adarsh K (02FE23BCI014) & KLE Technological University, Dr. MSSCET \\
\hline
\end{tabular}
\end{table}

\end{document}
